\documentclass{article}
\title{REPORT}
\author{Teenu Anand \& Geetha sri Koumudi}
\date{\today}
\begin{document}
\maketitle
\tableofcontents
\section{External Libraries}
 None
\section{Features}
\begin{enumerate}
    \item Secure SHA-256 authentication of players via Bash
    \item Three GUI games: $10 \times 10$,$8 \times 8$Othello,$7 \times 7$ Connect Four
    \item Displaying leaderboard after every game and data visualation using matplotlib
    \item General class design that is common for all games.so that we can add new games without changing overall thing.
    \item Interactive user interface(GUI) allows user to select games 
\end{enumerate}
\section{Design and Architecture}
\subsection{Architectural Layers}
\begin{itemize}
    \item Gatekeeper Layer (Bash): Manages \texttt{main.sh}. This handles file-stream operations (grep/awk) to validate users against \texttt{users.tsv} before launching the resource-heavy Python environment.
    \item Logic Layer (Python/NumPy): The core engine processes game rules. By representing boards as NumPy arrays, the system uses vectorized summation for O(1) or O(N) win-detection.
    \item Presentation Layer (Pygame): Maps internal states to visual pixels and manages the Finite State Machine (FSM) for transitions between menus and active play.
\end{itemize}
\section{Game implementations}
\subsection{Authentication system}
The main.sh handles the user authentication system ensuring only registered players can access the game
\begin{enumerate}
    \item it also makes sure every player has different usernames and store them in users.tsv.
    \item If a user exists,it triggers the login\_player.if not it offers registration for that player.
    \item It also avoids Player2 having same name as Player1 to avoid identity conflicts. 
    \item Upon successful authentication of both players it launches python game engine.
\end{enumerate}
\subsection{baseclass}
baseclass handles some features common to all games like
\begin{enumerate}
    \item stores player names(p1,p2) and whose turn(self.turn) is it.
    \item creates a gameboard in form of numpy array of size $n \times n$ initialised with empty cells
    \item A method to switch turns(switch\_turn()) between players.
    \item A method to save results of game in history.csv
\end{enumerate}
\subsection{tic-tac-toe}
This $10 \times 10$ board game where two players draw X and O alternatively whoever first gets five X in a row wins
\subsubsection{board setup}
\begin{enumerate}
    \item The board is initialised to be $10 \times 10$ grid using baseclass
    \item All places are set to be empty
\end{enumerate}
\subsubsection{makemove}
\begin{enumerate}
    \item Checks if cell is empty
    \item If valid,assigns the current player's mark to cell
    \item Returns True if the move is succesful,otherwise False
\end{enumerate}
\subsubsection{win detection}
\begin{enumerate}
    \item The board is converted into binary matrix where 1 represents current player's positions and 0 represent all other cells
    \item Function checks for 4 directions(horizontal,vertical,diagnol)
    \item if any segment sums for 5,a win is detected
    \item If there are no empty cells left on the board and no winner is found the game is declared draw
\end{enumerate}
\subsubsection{Rendering the board}
It handles all the graphical updates
\begin{enumerate}
    \item Draws the grid lines and Displays PLayers marks X for player 1 and O for player 2
    \item Displays the current game status(turn/winner/draw)
    \item Uses dynamic background colors based on player turns 
\end{enumerate}
\subsubsection{End screen}
\begin{enumerate}
    \item Displays final result
    \item Provides 2 options for user Play again(restarts the game) and Home(exists the game)
    \item Uses mouse detection for button interaction
\end{enumerate}
\subsubsection{Main Game loop}
The main function controls the execution flow
\begin{enumerate}
    \item Initialise the game window and fonts
    \item Runs a loop that \begin{enumerate}
        \item Handles user input(mouse clicks)
        \item converts screen coordinates to board positions
        \item Executes moves and checks for win/draw
        \item Updates the display continously
    \end{enumerate}
    \item After the game is completed directs to the end screen
\end{enumerate}
\subsection{Othello}
This 8X8 board game where 2 players take turns placing discs.when a move sandwitches opponent pieces,those pieces get flipped to the players pieces and the game ends when there are no moves are left
\subsubsection{board setup}
\begin{enumerate}
    \item The board is initialised as $8 \times 8$ grid
    \item The four central positions are filled with 2 discs for each player(standard starting configuration)
\end{enumerate}
\subsubsection{valid move detection}
\begin{enumerate}
    \item iterate through all emptycells
    \item for each cell check 8 possible directions(horizontal,vertical,diagnol)
    \item a move is valid if the adjacent cell is of opponent's piece and moving in that direction eventually reaches the current player's disc
\end{enumerate}
\subsubsection{making move}
\begin{enumerate}
    \item verifies if selectes move is valid
    \item Places the player's disc on board(changing numpy values to 1 or 2 based on whose piece is it)
    \item For each direction it collects opponents dics in a straight line and turns them into player's disc
    \item switches after a move and if the player has no valid moves then the turn is skipped
\end{enumerate}
\subsubsection{drawing}
\begin{enumerate}
    \item Uses pygame to\begin{enumerate}
        \item Draw the game board
        \item Display discs(black/white)
        \item Hoighlight valid moves
    \end{enumerate}
    \item Displays player scores and Current turn
\end{enumerate}
\subsubsection{Checking win}
\begin{enumerate}
    \item Game ends when both players dont have any move left
    \item Higher score count player is declared as win
\end{enumerate}
\subsection{Connect4}
This $7 \times 7$ board game drops discs from the top and whoever gets contionous 4 in any direction gets to win
\subsubsection{Board configuration}
\begin{enumerate}
    \item The board is initialised by $7 \times 7$ grid
    \item All are initialsed to 0
    \item Discs are dropped vertically and occupy the lowest positons available in column 
\end{enumerate}
\subsubsection{Move validation}
\begin{enumerate}
    \item Checks if the topmost cell of column is empty
    \item Ensures disc can be dropped in the selected column
\end{enumerate}
\subsubsection{Finding the next Open row}
\begin{enumerate}
    \item Iterates from the bottom of the column upwards
    \item Returns the first available empty row
\end{enumerate}
\subsubsection{Dropping a disc}
\begin{enumerate}
    \item Uses nextopenrow() to determine placement
    \item Assigns the current player's value to the identified cell
\end{enumerate}
\subsubsection{win detection}
\begin{enumerate}
    \item The board is converted into a binary matrix where:
    \begin{enumerate}
        \item 1 represents the current player's discs
        \item 0 represents all other cells
    \end{enumerate}
    \item The function checks for 4 directions(horizontal,vertical,main diagnol,anti diagnol)
    \item Uses numpy slicing and summation to efficiently detect sequences of length 4
    \item Game ends in draw if all the board cells are filled and no win condition was satisfied
\end{enumerate}
\subsubsection{Rendering the game}
\begin{enumerate}
    \item Drawing the grid using blue rectangles
    \item Rendering empty slots as black circles
    \item Displaying player discs:
    \begin{enumerate}
        \item Red for Player 1
        \item Yellow for Player 2
    \end{enumerate}
    \item Positioning discs correctly to simulate gravity (bottom-up filling)
\end{enumerate}
\subsubsection{Main game loop}
\begin{enumerate}
    \item Initializes the game window and rendering parameters
    \item Continously takes User's input(Mouse clicks)
    \item Checks for Win condition
    \item Displays the winner or draw message when the game ends
\end{enumerate}
\subsubsection{leaderboard appearing}
\begin{enumerate}
    \item Updates the users.tsv using the save result function and gives the leaderboard
\end{enumerate}
\subsection{leaderboard}
Tracks and ranks player performance across games
\begin{enumerate}
    \item History.csv contains details of all the games(stored in form of comma seperated winner,loser,game,date)
    \item while showing leaderboard to interface it is sorted based on user's demand(win,lose,$\frac{win}{loss}$ratio)
    
\end{enumerate}
\section{Project Structure}
\begin{verbatim}
Mini_Hub_GTA1/
|-- main.sh              
|-- game.py              
|-- leaderboard.sh       
|-- history.csv          
|-- users.tsv            
|-- home.png            
|-- tictactoe.png        
|-- othello.png        
|-- connect4.png       
+-- games/             
    |-- base_class.py  
    |-- tictactoe.py    
    |-- othello.py   
    +-- connect4.py    
\end{verbatim}
\section{Installation and setup}
\subsection{Prerequisites}
\begin{enumerate}
    \item Python 10.x installed on the system
    \item Required libraries:
    \begin{enumerate}
        \item \texttt{pygame}
        \item \texttt{numpy}
    \end{enumerate}
\end{enumerate}
\subsection{Installation}
\begin{enumerate}
    \item Clone or download the project repository
    \item Navigate to the project directory
    \item Install dependencies using:
\begin{verbatim}
pip install pygame numpy
\end{verbatim}
\end{enumerate}
\subsection{running the game}
\begin{enumerate}
    \item The system is launched using the shell script
    \begin{verbatim}
        bash main.sh
    \end{verbatim}
    \item It prompts for usernames and passwords and then opens game interface.players interact with mouse 
\end{enumerate}
\section{Hurdles}
\begin{enumerate}
    \item Implementing win detection across multiple directions (horizontal, vertical, diagonal) .
    \item merge took lot of tome to sort.due to us writing independently use of different variables and somrtimes different logics
    \item coordinate transformations between screen space and board indices was error-prone.
    \item Ensuring correct turn switching
\end{enumerate}
\section{How will we use 10 hours extra time}
\begin{enumerate}
    \item Would have tried to include a extra game like chess
    \item Could have added extra design effects and made interface to look aesthetic
    \item May be tried having a bot as an opponent 
    \item Could have made dyanmic board sizes
\end{enumerate}
\section{Bibliography}
\begin{enumerate}
 \item NumPy Developers, NumPy Documentation.  
    Available at: \texttt{https://numpy.org/doc/}

 \item Pygame Community, Pygame Documentation.  
    Available at: \texttt{https://www.pygame.org/docs/}
\end{enumerate}

\end{document}